%% notes-tex/kg-build-system/sections/2-build-host.tex

\section{The build host\secstatus{todo}}
\label{sec:build}

\notesec{What the build needs}
The full build is \file{database/slurm/scripts/gilahyper_live_rebuild-slurm_docker.slurm}.
It reads the development LMDBs read-only, checks that each one's build manifest
is fresh under the schema surface the container will install, generates the
BioCypher CSVs in a build container pinned to the build commit, imports them with
\texttt{neo4j-admin database import} into a new store beside the served one,
validates it, swaps it under the serving container, and stamps it: the release id
\texttt{<build date>-<commit[:8]>}, a \texttt{KgRelease} node carrying the version,
the package version and tag, and one content hash per dataset, the aliases
\texttt{latest} and \texttt{pinned}, and the committed snapshot. The old store
stays beside the new one as \texttt{data.superseded.<timestamp>} for rollback.
Table~\ref{tab:hostneeds} lists what that script assumes of its host, which is the
checklist for moving it.

%% SOURCE: hand-authored from the slurm script and the measurements of job 3297
%% (notes/database.slurm.scripts.gilahyper_live_rebuild-slurm_docker.md).
\begin{table}[htbp]
\centering
\footnotesize
\caption{What the build host provides today, and whether a new host must match it.}
\label{tab:hostneeds}
\begin{tabular}{>{\raggedright\arraybackslash}p{0.24\textwidth}>{\raggedright\arraybackslash}p{0.36\textwidth}>{\raggedright\arraybackslash}p{0.3\textwidth}}
\toprule
Need & GilaHyper today & When the build moves \\
\midrule
Development LMDB tree & \texttt{\$DATA\_ROOT/}\file{data/torchcell/<dataset>/} on \texttt{/scratch}, one store per mapped dataset, each with a fresh \file{build_manifest.json} & must exist on the new host; the loaders rebuild it from the raw mirror, or it is copied \\
Scheduler and container runtime & slurm with \texttt{--cpus-per-task=64 --mem=256G}; docker with the \texttt{tc-neo4j} image & any scheduler that runs one job with that budget; docker or an equivalent that honors the same mounts \\
Store disk & \texttt{/db}, 7.3 TB NVMe; the import writes the new store beside the served one & local disk, not NFS; the 3.0 store is 151 GB restored, and a build holds two \\
Archive disk & \texttt{/bulk}, 24 TB RAID, mounted into the container at \texttt{/kg-releases}; the online backup's work directory holds the whole store uncompressed & any local disk the container can mount; the work directory needs the store's size again \\
Fingerprint checkout & the primary checkout at the build commit, which must be a tagged commit & the same; the tag gate is in the script \\
Serving container & \texttt{tc-neo4j-readonly}, swapped in place with minutes of downtime & a build host need not serve; a build-only host ends at the artifact and ships it \\
\bottomrule
\end{tabular}
\end{table}

\notesec{What the last build measured}
Job 3297 built release \texttt{2026.10.06-4b293d34} (KG 3.0): 51 datasets,
100,127,515 nodes, CSV generation 3 h 09 min at 64 CPUs and 256 GB, the served
store 151 GB on disk once restored. The online backup of that store is
11,175,064,969 bytes, and its copy to Taiga took 2 min 27 s at 111 MB/s. A restore
of that artifact into a second database on the same DBMS, with serving
uninterrupted, took 10 min 13 s (jobs 3338 and 3339).

\notesec{Build-only versus build-and-serve}
GilaHyper does both, which is why the script ends with a swap. A build-only host
would stop after validation, run \texttt{kg\_release.sh backup}, \texttt{archive}
and \texttt{ship}, and let every serving host, GilaHyper included, pick the
release up with \texttt{deploy}. Nothing in the pairing depends on which of the
two shapes the build host has: the snapshot, the node and the tag are written
either way.
